\chapter*{Introduction générale}
\addcontentsline{toc}{chapter}{Introduction générale}
\markboth{Introduction générale}{}

\section*{Contexte}

Le Maroc a fait de l'intelligence artificielle souveraine une priorité nationale, portée
notamment par la feuille de route «~Maroc Digital 2030~», l'initiative «~AI Made in Morocco~» et
le réseau des instituts Al Jazari. L'un des principes directeurs de cette stratégie est la
\textbf{maîtrise de la donnée} : une information sensible produite au Maroc doit pouvoir être
traitée sur le territoire, sans dépendance critique vis-à-vis d'infrastructures étrangères.

Ce principe prend un relief particulier dans le domaine juridique. Une question posée par un
citoyen sur ses droits --- son licenciement, son congé de maternité, son salaire --- n'est pas
une requête anodine : c'est une information intime, souvent formulée dans un moment de
vulnérabilité. La confier à un service en ligne étranger revient à la faire sortir du territoire,
avec les conséquences que cela emporte en matière de confidentialité et de dépendance
technologique.

\section*{Problématique}

L'accès du citoyen à l'information juridique reste aujourd'hui difficile. Les textes --- Code du
travail, Code de la famille, Bulletin Officiel --- sont publics, mais rédigés dans une langue
dense, structurés en centaines d'articles se renvoyant les uns aux autres, et parfois disponibles
uniquement en arabe. Il n'existe pas d'outil largement accessible permettant de répondre à une
question simple, posée en langage courant, tout en garantissant que la réponse provient bien du
texte officiel.

Les assistants conversationnels généralistes ne comblent pas ce manque, pour deux raisons
distinctes. D'une part, ils connaissent mal le droit marocain et peuvent \textbf{inventer des
références inexistantes} : un numéro d'article plausible mais faux est plus dangereux qu'un aveu
d'ignorance, parce que rien dans sa forme ne le distingue d'une référence exacte. D'autre part,
ils font transiter les questions des citoyens vers des serveurs étrangers, en contradiction
directe avec l'exigence de souveraineté rappelée plus haut.

La question à laquelle ce stage s'attaque peut donc se formuler ainsi : \emph{comment construire
un système qui répond à une question juridique en langage naturel, en garantissant que chaque
affirmation est rattachée à un texte officiel vérifiable, et sans qu'aucune donnée ne quitte la
machine~?}

\section*{Contribution}

Ce stage, effectué au sein de Pulsaride Solutions du 1\textsuperscript{er}~juillet au
31~août~2026, a consisté à concevoir, réaliser et évaluer un assistant de recherche juridique
souverain fondé sur l'architecture RAG (\anglais{Retrieval-Augmented Generation}), appliqué au
Code du travail marocain (Loi n\textsuperscript{o}~65-99). Le système livré :

\begin{itemize}
  \item s'appuie sur un corpus de \textbf{\NbArticles{} articles} construit depuis le PDF
        officiel, reproductible à l'identique ;
  \item retrouve les articles pertinents en combinant recherche sémantique et recherche
        lexicale ;
  \item génère une réponse fondée \emph{uniquement} sur ces articles, en citant systématiquement
        leurs numéros ;
  \item \textbf{s'abstient} explicitement lorsque l'information demandée ne figure pas dans le
        corpus ;
  \item s'exécute \textbf{entièrement en local}, sans aucun appel à une interface de
        programmation externe.
\end{itemize}

L'apport ne se limite pas au prototype. Le stage a produit un \textbf{banc de mesure rejouable}
qui a permis d'établir des constats que ni l'intuition ni des tests ponctuels n'auraient
révélés --- au premier rang desquels une hallucination inter-codes que les tests manuels n'avaient
jamais déclenchée.

Ce banc a lui-même fait l'objet d'un audit, et c'est peut-être le résultat le plus instructif du
stage : son jeu de questions initial n'exerçait pas plusieurs comportements du système, et son
élargissement a \emph{inversé} deux conclusions que l'on croyait établies, sans qu'une seule ligne
du moteur ait changé. Ces résultats, et la manière dont ils ont été obtenus, occupent une place
centrale dans ce rapport.

\section*{Organisation du rapport}

Le rapport est structuré en cinq chapitres. Le
chapitre~\ref{chap:contexte} présente le cadre du stage : l'organisme d'accueil, les objectifs,
le périmètre volontairement borné et l'organisation du travail. Le chapitre~\ref{chap:fondements}
expose les fondements techniques mobilisés, et en particulier la raison pour laquelle un grand
modèle de langage seul ne peut pas répondre de façon fiable à une question juridique. Le
chapitre~\ref{chap:conception} traduit ces principes en une architecture concrète et justifie les
choix technologiques, y compris ceux qui restent discutables. Le chapitre~\ref{chap:realisation}
décrit la réalisation, de la constitution du corpus aux garde-fous de génération, en accordant
une place explicite aux difficultés rencontrées. Le chapitre~\ref{chap:evaluation} présente le
protocole d'évaluation et l'ensemble des résultats mesurés, limites comprises. Une conclusion
générale dresse le bilan du stage et ses perspectives.
